% !TEX program = xelatex
% ============================================================
%  سجل أصحاب المصلحة — Stakeholder Register Template
%  نموذج لتسجيل أصحاب المصلحة في المشروع وتحليل مصالحهم وتأثيرهم
%  Compile with: xelatex main.tex (run twice)
% ============================================================
%  Features:
%    - Stakeholder table: Name | Role | Organization | Interest |
%      Influence (High/Med/Low) | Engagement strategy | Contact
%    - 7 sample rows
%    - Influence-interest summary section
% ============================================================

\documentclass[11pt,a4paper]{article}

\usepackage[a4paper,margin=2.5cm,top=2.5cm,bottom=2.5cm,headheight=16pt]{geometry}
\usepackage{graphicx}
\usepackage{array}
\usepackage{booktabs}
\usepackage{tabularx}
\usepackage{multirow}
\usepackage{enumitem}
\usepackage{float}
\usepackage{titlesec}
\usepackage{fancyhdr}
\usepackage{setspace}
\usepackage{xcolor}
\usepackage{amsmath}

\usepackage{bidi}
\usepackage[no-math]{fontspec}
\usepackage[quiet]{polyglossia}

\setmainfont[Scale=1]{Amiri-Regular.ttf}
\newfontfamily\arabicfont[Script=Arabic,Scale=0.95]{Amiri-Regular.ttf}
\newfontfamily\englishfont[Scale=0.95]{TeX Gyre Termes}

\setdefaultlanguage[calendar=gregorian,hijricorrection=1]{arabic}
\setotherlanguage[variant=british]{english}

\usepackage[breaklinks,hidelinks]{hyperref}

\addto\captionsarabic{%
  \renewcommand{\contentsname}{الفهرس}%
  \renewcommand{\figurename}{الشكل}%
  \renewcommand{\tablename}{الجدول}%
}

\titleformat{\section}{\Large\bfseries}{\thesection}{0.7em}{}
\titlespacing*{\section}{0pt}{1.4ex plus 0.4ex}{0.9ex plus 0.3ex}

\pagestyle{fancy}
\fancyhf{}
\fancyhead[R]{\small\itshape سجل أصحاب المصلحة}
\fancyhead[L]{\small اسم المشروع}
\fancyfoot[C]{\small\thepage}
\renewcommand{\headrulewidth}{0.3pt}

\setlist[itemize]{topsep=0.3em, itemsep=0.15em, parsep=0pt}
\setlist[enumerate]{topsep=0.3em, itemsep=0.15em, parsep=0pt}

\newcolumntype{L}[1]{>{\raggedright\arraybackslash}p{#1}}
\newcolumntype{C}[1]{>{\centering\arraybackslash}p{#1}}

\setstretch{1.3}

\begin{document}

% ============================================================
% TITLE BLOCK
% ============================================================
\begin{center}
{\LARGE\bfseries سجل أصحاب المصلحة}\\[0.3em]
{\large Stakeholder Register}\\[1em]
\end{center}

% ============================================================
% PROJECT INFO
% ============================================================
% TODO: املأ بيانات المشروع.
\begin{table}[H]
\centering
\renewcommand{\arraystretch}{1.4}
\begin{tabularx}{\textwidth}{L{3.5cm} X L{3.5cm} X}
\toprule
\textbf{اسم المشروع:} & --- & \textbf{مدير المشروع:} & --- \\
\textbf{اسم الشركة:} & --- & \textbf{تاريخ التحديث:} & --- \\
\bottomrule
\end{tabularx}
\end{table}

\vspace{1em}

% ============================================================
% 1. STAKEHOLDER REGISTER TABLE
% ============================================================
\section{سجل أصحاب المصلحة}
\label{sec:register}

% TODO: عدّل صفوف أصحاب المصلحة حسب مشروعك.
\begin{table}[H]
\centering
\caption{سجل أصحاب المصلحة في المشروع}
\renewcommand{\arraystretch}{1.4}
\small
\begin{tabularx}{\textwidth}{L{2.5cm} L{2.5cm} L{2.5cm} L{3cm} C{1.3cm} L{3cm} L{2.5cm}}
\toprule
\textbf{الاسم} & \textbf{الدور} & \textbf{الجهة} & \textbf{المصلحة} & \textbf{التأثير} & \textbf{استراتيجية الإشراك} & \textbf{التواصل} \\
\midrule
--- & الراعي الرئيسي & الإدارة العليا & تمويل المشروع ونجاحه & عالي & تقارير شهرية واجتماعات ربع سنوية & بريد إلكتروني \\
--- & مدير المشروع & إدارة المشاريع & تنفيذ المشروع وتسليمه & عالي & اجتماعات أسبوعية وتقارير يومية & بريد / هاتف \\
--- & قائد الفريق التقني & الفريق التقني & الجدوى التقنية والتنفيذ & متوسط & اجتماعات أسبوعية ومراجعات فنية & بريد إلكتروني \\
--- & المسؤول المالي & الإدارة المالية & الالتزام بالميزانية & متوسط & تقارير شهرية ومراجعة المصروفات & بريد إلكتروني \\
--- & ممثل المستخدمين & المستفيدون النهائيون & سهولة الاستخدام والوظائف & منخفض & استبيانات واجتماعات دورية & بريد / اجتماعات \\
--- & مدير الجودة & إدارة الجودة & جودة المخرجات & متوسط & مراجعات قبل التسليم وتدقيق دوري & بريد إلكتروني \\
--- & المورد الخارجي & شركة المورد & الالتزامات التعاقدية & منخفض & اجتماعات شهرية وتقارير تسليم & بريد / هاتف \\
\bottomrule
\end{tabularx}
\end{table}

% ============================================================
% 2. INFLUENCE-INTEREST SUMMARY
% ============================================================
\section{تصنيف التأثير والمصلحة}
\label{sec:classification}

% TODO: حدّث تصنيف أصحاب المصلحة حسب التحليل الفعلي.
\begin{table}[H]
\centering
\caption{تصنيف أصحاب المصلحة حسب التأثير والمصلحة}
\renewcommand{\arraystretch}{1.4}
\begin{tabularx}{\textwidth}{L{4cm} C{3cm} C{3cm} X}
\toprule
\textbf{الفئة} & \textbf{التأثير} & \textbf{المصلحة} & \textbf{استراتيجية الإدارة} \\
\midrule
الراعي الرئيسي & عالي & عالي & إدارة عن قرب --- إبقائه راضياً ومشتركاً \\
مدير المشروع & عالي & عالي & إدارة عن قرب --- شريك رئيسي في التنفيذ \\
قائد الفريق التقني & متوسط & عالي & إبقائه مطّلعاً --- استشار في القرارات التقنية \\
المسؤول المالي & متوسط & متوسط & إبقائه راضياً --- تقارير دورية عن الميزانية \\
ممثل المستخدمين & منخفض & عالي & إبقيه مطّلعاً --- استشر في المتطلبات \\
مدير الجودة & متوسط & متوسط & إبقيه راضياً --- شاركه في مراجعات الجودة \\
المورد الخارجي & منخفض & منخفض & جهد معتدل --- رصد الأداء التعاقدي \\
\bottomrule
\end{tabularx}
\end{table}

% ============================================================
% 3. ENGAGEMENT NOTES
% ============================================================
\section{ملاحظات الإشراك}
\label{sec:notes}

% TODO: أضف ملاحظات حول استراتيجية إشراك أصحاب المصلحة.
\begin{itemize}
    \item الراعي الرئيسي ومدير المشروع هما أصحاب التأثير والمصلحة الأعلى --- يتطلبان إشراكاً مستمراً.
    \item ممثل المستخدمين ذو تأثير منخفض لكن مصلحة عالية --- يجب إبقاؤه مطّلعاً لضمان قبول المنتج.
    \item المورد الخارجي يتطلب رصداً منتظماً لضمان الالتزام بمواعيد التسليم.
\end{itemize}

\end{document}
